\clearpage
\begingroup
\fancyhead[L]{\small THEORIE DER BEWEGUNG DES MONDES.}
\renewcommand{\headrulewidth}{0pt}
\clearpage
\SourceLeaf{617}
\begin{center}
\vspace*{.15\textheight}
{\Large\bfseries THEORIE DER BEWEGUNG DES MONDES.\par}
\end{center}
\vspace{.28\textheight}
\par\noindent\hfill 77\textsuperscript{*}\par

\clearpage
\SourceLeaf{619}
\begin{center}
{\large NACHLASS.\par}
\bigskip
{\Large\bfseries THEORIE DER BEWEGUNG DES MONDES.\par}
\bigskip
ERSTER ABSCHNITT.
\medskip

\textit{Fundamentalgleichungen für die Bewegung des Mondes.}
\bigskip

{\large Vorerinnerung.\par}
\end{center}

\par\noindent\hspace*{\parindent}Von den Kräften, die auf die Bewegung des Mondes Einfluss haben, betrachten
wir hier nur die Anziehungskräfte der Erde, der Sonne und des Mondes, welche
wir als in den Schwerpunkten dieser Weltkörper vereinigt ansehen. In der
That sind die Wirkungen der übrigen Planeten so ausserordentlich klein, dass
sie auch den allerfeinsten Beobachtungen gänzlich entgehen; die sphäroidische
Gestalt der Erde, wodurch die letztere Voraussetzung modificirt wird, bringt
auch nur ein Paar Gleichungen von wenigen Secunden hervor. Zudem ist es
nicht schwer, auch von diesen Umständen Rechnung zu tragen, sobald die
Bewegung des Mondes einmal ohne dieselben erörtert ist.

\par\noindent\hspace*{\parindent}Wir werden ausserdem die Bewegung der Erde um die Sonne als elliptisch
annehmen, und die Störungen derselben, durch den Mond und die Planeten, bei
Seite setzen. Wir betrachten mithin auch die Bahn der Erde um die Sonne als
in einer unveränderlichen (d. i. sich selbst beständig parallel bleibenden)
Ebene, und vernachlässigen sowohl das allmähliche Déplacement
\SourceLeaf{620}dieser Ebene selbst, welches die Einwirkung der Planeten hervorbringt,
als die periodischen Ausbeugungen der Erde von derselben, die eine Folge der
Perturbationen der Planeten und des Mondes zugleich sind. Denn alle diese
Umstände erzeugen in der Bewegung des Mondes um die Erde Modificationen, die
für sich unmerklich sind, und ganz füglich einer eignen Untersuchung
aufgespart werden können.

\par\noindent\hspace*{\parindent}Zur Aufsuchung der Fundamentalgleichungen beziehen wir zuvörderst die
Örter des Mondes, der Erde und der Sonne im Weltraum auf drei unveränderliche
Ebenen, die sich in einem ganz willkürlichen oder vielmehr unbestimmt
bleibenden Punkte des Weltraumes unter rechten Winkeln schneiden. Die eine
dieser Ebenen legen wir der Ekliptik parallel und bezeichnen die Entfernungen
des Mondes, der Erde und der Sonne von derselben zu jeder Zeit \(t\) durch
\(z,z',z''\) resp.; diese Abstände betrachten wir als positiv auf der Seite, wo
der Nordpol der Ekliptik liegt, als negativ auf der entgegengesetzten. Die
zweite Ebene legen wir senkrecht auf die vorige, so, dass die
Durchschnittslinie durch den Punkt der Himmelskugel geht, wo der Frühlingspunkt
zu irgend einer bestimmten Zeit \(T\) ist; die Abstände des Mondes, der Erde und
der Sonne von dieser Ebene, positiv auf der Seite genommen, wo der
Sommersonnenwendepunkt liegt, bezeichnen wir mit \(y,y',y''\). Endlich wird die
dritte Ebene senkrecht auf die beiden vorigen sein, und ihre Durchschnittslinie
mit der Erstern durch den Sommersonnenwendepunkt gehen; die Abstände des
Mondes, der Erde und der Sonne von derselben sehen wir auf der Seite des
Frühlingsnachtgleichepunkts als positiv, auf der entgegengesetzten als negativ
an und heissen sie \(x,x',x''\).

\par\noindent\hspace*{\parindent}Nennen wir nunmehr die Abstände der Erde von der Sonne, des Mondes von der
Sonne, des Mondes von der Erde resp. \(R,R',R''\), und die Massen des Mondes, der
Erde und der Sonne resp. \(m,m',m''\), so geben die bekannten Grundsätze der
Dynamik folgende sechs Fundamentalgleichungen
\[
\begin{aligned}
0&=\frac{\mathrm{dd}x}{\mathrm{d}t^2}
  +\frac{m'(x-x')}{R''^3}+\frac{m''(x-x'')}{R'^3}\\[4pt]
0&=\frac{\mathrm{dd}y}{\mathrm{d}t^2}
  +\frac{m'(y-y')}{R''^3}+\frac{m''(x-x'')}{R'^3}\\[4pt]
0&=\frac{\mathrm{dd}z}{\mathrm{d}t^2}
  +\frac{m'(z-z')}{R''^3}+\frac{m''(z-z'')}{R'^3}
\end{aligned}
\]
\SourceLeaf{621}
\[
\begin{aligned}
0&=\frac{\mathrm{dd}x'}{\mathrm{d}t^2}
  +\frac{m(x'-x)}{R''^3}+\frac{m''(x'-x'')}{R^3}\\[4pt]
0&=\frac{\mathrm{dd}y'}{\mathrm{d}t^2}
  +\frac{m(y'-y)}{R''^3}+\frac{m''(y'-y'')}{R^3}\\[4pt]
0&=\frac{\mathrm{dd}z'}{\mathrm{d}t^2}
  +\frac{m(z'-z)}{R''^3}+\frac{m''(z'-z'')}{R^3}
  \mathpunct{\raisebox{.4ex}{.}}
\end{aligned}
\]
\clearpage
\endgroup
